% ======================================================================
% INSTRUCTIONS WHEN EDITING THIS PROJECT
% ======================================================================
%
% 1. Do not guess labels or numbers.
%    Always use \ref{...} exactly as requested. Never invent labels.
%
% 2. When asked to insert a reference, write it literally as
%       Thought Experiment~\ref{te:lorentz}
%    or whichever label is given, without adding chapter, section, or numbers.
%
% 3. Use exact hierarchy:
%       \chapter{...}
%       \section{...}
%       \subsection{...}
%       \subsubsection{...}
%    Never invent section names or change levels.
%
% 4. Single, clean LaTeX block output only.
%    No prose, no explanation, no commentary outside the code block.
%
% 5. ASCII only in body text. No Unicode punctuation.
%
% 6. BibTeX citation keys are lastnameYYYY and must be sorted alphabetically
%    inside each citation command.
%
% 7. SOFT MARGINS OF 100 characters for all generated prose.
%
% 8. Never write explanations, apologies, or meta-text unless the user asks.
%    Do not self-justify, speculate, question instructions, or add commentary.
%    The output should assume the user already knows what they want and expects
%    exact compliance, not discussion.
%
% 9. Add labels to all equations, definitions, sections, chapters, and sidebars:
%      \label{eq:[SHORT-NAME-WITH-DASHES]}
%      \label{def:[SHORT-NAME-WITH-DASHES]}
%      \label{se:[SHORT-NAME-WITH-DASHES]}
%      \label{chap:[SHORT-NAME-WITH-DASHES]}
%      \label{sb:[SHORT-NAME-WITH-DASHES]}
%
%10. Lean code is set in lean environments, not generic verbatim.
%    Prose references to Lean names use \lean{name}.
%
%11. Do not gloss words.
% ======================================================================
% END OF INSTRUCTIONS
% ======================================================================


\newtheorem{theorem}{Theorem}
\newtheorem{axiom}{Axiom}
\newtheorem{remark}{Remark}
\newtheorem{definition}{Definition}
\newtheorem{proposition}{Proposition}
\newtheorem{corollary}{Corollary}
\newtheorem{law}{Law}

% Lean inline and display
\newcommand{\lean}[1]{\texttt{#1}}
\newenvironment{leancode}{\begin{quote}\ttfamily\small}{\end{quote}}

% Sidebar / callout environment for non-technical readers
\newenvironment{sidebar}[1]{%
  \begin{center}
  \begin{minipage}{0.88\linewidth}
  \hrule\vspace{0.5em}
  \textbf{#1}\par\smallskip
}{%
  \vspace{0.5em}\hrule
  \end{minipage}
  \end{center}
  \medskip
}

% Episode reference
\newcommand{\episode}[1]{\textit{Episode~#1}}

% N.B. callout
\newcommand{\NB}[1]{%
  \par\noindent\textbf{N.B.---}#1\hfill$\square$\par
}

% Ledger symbol
\newcommand{\Ledger}{\mathcal{L}}

% Heartbeat annotation (the bullshit meter)
\newcommand{\heartbeat}[1]{\marginpar{\tiny $\heartsuit \approx #1$}}

\makeatletter
\renewcommand{\l@chapter}{\@dottedtocline{0}{0em}{1.5em}}
\makeatother

\hyphenation{treat-ed}
